\documentclass[11pt,a4paper]{article}
\usepackage[margin=2.5cm]{geometry}
\usepackage{amsmath,amssymb}
\usepackage{graphicx}
\usepackage{booktabs}
\usepackage{hyperref}
\usepackage[numbers]{natbib}

\title{Forecasting the S\&P 500 One Day Ahead:\\ARIMA, a Vanilla RNN and an LSTM Against a Random-Walk Baseline}
\author{Name 1 (ID) \and Name 2 (ID) \and Name 3 (ID)}
\date{COSC 621, Data Science and Generative AI \\ \today}

\begin{document}
\maketitle

\begin{abstract}
One paragraph: the question (can ARIMA, a vanilla RNN and an LSTM forecast the next day's S\&P 500
close better than a random walk?), the data (source, period, sample size), the design (chronological
split, rolling one-day forecasts, errors in index points), the headline result and its interpretation.
\end{abstract}

\section{Introduction}
Why forecasting an index is hard, what a random walk implies for forecasting, what this report tests
and how the rest is organised.

\section{Data and analysis}
\subsection{Data collection and preparation}
Source, series identifier, access date, period and its justification, sample size, checks for order,
duplicates, missing observations, market closures versus missing data, outliers, transformations.
\subsection{Time series analysis}
Trend, seasonality (how detected, evidence found or absent), stationarity (ADF and KPSS with their
hypotheses and limitations), ACF and PACF, the random walk defined formally with its innovation
assumptions, and whether the evidence is consistent with it.

\section{Methodology}
\subsection{Forecasting design}
Target (next-day close; conversion from log returns), split dates and their justification, rolling
forecast and refitting procedure, preprocessing fitted on training data only.
\subsection{Random-walk baseline}
\subsection{ARIMA}
Mechanism and equations; p, d, q justified by stationarity checks, ACF/PACF, information criteria and
validation; residual inspection.
\subsection{Vanilla RNN}
Recurrence and memory; lookback window, inputs, layers, hidden units, loss, optimizer, learning rate,
batch size, epochs, overfitting controls; configuration table.
\subsection{LSTM}
Gates (forget, input, output) and why they help; same justification and configuration table.

\section{Results and discussion}
Error comparison table (MAE and RMSE in index points, identical test dates), actual-versus-predicted
plots, residual patterns, overfitting, computational cost, whether any improvement is meaningful,
limitations.

\section{Conclusion}

\bibliographystyle{plainnat}
\bibliography{refs}
\end{document}
